% Part: model-theory
% Chapter: interpolation
% Section: definability

\documentclass[../../../include/open-logic-section]{subfiles}

\begin{document}

\olfileid[fr]{mod}{int}{def}

\olsection{Le théorème de définissabilité}

Le théorème de définissabilité de Beth
est une application importante du
théorème d'interpolation. Pour définir
une relation~$R$ à $n$ places,
on peut donner une !!{formula}~$!C$
à $n$ !!{variable}s libres qui ne
fait pas intervenir~$R$. C'est une
définition \emph{explicite} de~$R$
au moyen de~$!C$. De même, une
théorie~$\Sigma(P)$ dont le langage
contient le !!{predicate}~$P$ à
$n$ places définit explicitement~$P$
si elle contient, ou du moins entraîne,
une définition formalisée :
\[
\Sigma(P)\Entails\lforall[x_1][\dots
  \lforall[x_n][(\Atom{P}{x_1,\dots,x_n}
  \liff !C(x_1,\dots,x_n))]].
\]
Mais on peut aussi déterminer
entièrement une relation sans en
donner de formule explicite : il
se peut que, dans tout modèle de
la théorie, l'interprétation de~$P$
soit fixée par celle des autres
symboles du langage. On dit alors
que la théorie définit~$P$
\emph{implicitement}. Le théorème
de Beth affirme qu'une telle
définition donne toujours lieu
à une définition explicite.

\begin{defn}
Soit $\Lang{L}$ un !!{language}
ne contenant pas le !!{predicate}~$P$.
Un ensemble $\Sigma(P)$ d'%
!!{sentence}s de $\Lang{L}\cup\{P\}$
\emph{définit explicitement}~$P$
si et seulement s'il existe une
!!{formula}~$!C(x_1,\dots,x_n)$
de~$\Lang{L}$, dont les seules variables éventuellement libres sont $x_1,\dots,x_n$, telle que
\[
\Sigma(P)\Entails\lforall[x_1][\dots
  \lforall[x_n][(\Atom{P}{x_1,\dots,x_n}
  \liff !C(x_1,\dots,x_n))]].
\]
\end{defn}

\begin{defn}
Soit $\Lang{L}$ un !!{language}
ne contenant ni~$P$ ni~$P'$.
Un ensemble $\Sigma(P)$ d'%
!!{sentence}s de $\Lang{L}\cup\{P\}$
\emph{définit implicitement}~$P$
si et seulement si
\[
\Sigma(P)\cup\Sigma(P')\Entails
  \lforall[x_1][\dots
  \lforall[x_n][(\Atom{P}{x_1,\dots,x_n}
  \liff\Atom{P'}{x_1,\dots,x_n})]],
\]
où $\Sigma(P')$ s'obtient en
remplaçant partout~$P$ par~$P'$
dans $\Sigma(P)$.
\end{defn}

Autrement dit, soient une
$\Lang{L}$-structure $\Struct{M}$
et deux relations
$R,R'\subseteq\Domain{M}^n$.
Si l'expansion de~$\Struct{M}$
qui interprète~$P$ par~$R$
satisfait $\Sigma(P)$, et celle
qui interprète~$P'$ par~$R'$
satisfait $\Sigma(P')$, alors
$R=R'$. La définition implicite
exprime donc l'unicité de
l'interprétation compatible avec
la théorie pour une même structure
du langage de base.

\begin{thm}[Théorème de définissabilité de Beth]
Un ensemble $\Sigma(P)$ d'%
!!{sentence}s de $\Lang{L}\cup\{P\}$
définit implicitement~$P$ si et
seulement s'il le définit
explicitement.\footnote{Précision éditoriale : l'énoncé source dit ici « formules », alors que les deux définitions précédentes portent sur des ensembles d'énoncés fermés. Nous conservons ce cadre dans le théorème et sa preuve.}
\end{thm}

\begin{proof}
Supposons d'abord que $\Sigma(P)$
définisse explicitement~$P$ par
$!C(x_1,\dots,x_n)$. Le remplacement
uniforme de~$P$ par~$P'$ donne
\begin{align*}
  \Sigma(P)&\Entails
    \lforall[x_1][\dots\lforall[x_n]
    [(\Atom{P}{x_1,\dots,x_n}
    \liff !C(x_1,\dots,x_n))]],\\
  \Sigma(P')&\Entails
    \lforall[x_1][\dots\lforall[x_n]
    [(\Atom{P'}{x_1,\dots,x_n}
    \liff !C(x_1,\dots,x_n))]].
\end{align*}
Il en résulte que $\Sigma(P)$
définit implicitement~$P$.

Réciproquement, supposons la
définition implicite. Ajoutons au
langage des constantes nouvelles deux à deux distinctes
$c_1,\dots,c_n$. Par la définition,
\[
\Sigma(P)\cup\Sigma(P')\Entails
  \Atom{P}{c_1,\dots,c_n}
  \lif\Atom{P'}{c_1,\dots,c_n}.
\]
La compacité donne des ensembles
finis $\Delta_0\subseteq\Sigma(P)$
et $\Delta_1\subseteq\Sigma(P')$
qui entraînent déjà cette
implication. Remplaçons~$P'$ par~$P$
dans les !!{sentence}s de $\Delta_1$,
et soit~$!D(P)$ la conjonction
de ces !!{sentence}s et de ceux
de~$\Delta_0$. Une conjonction
vide désigne ici une tautologie.
Écrivons $!D(P')$ pour le résultat
du remplacement uniforme de~$P$
par~$P'$ dans~$!D(P)$. Alors
\[
!D(P)\land !D(P')\Entails
  \Atom{P}{c_1,\dots,c_n}
  \lif\Atom{P'}{c_1,\dots,c_n}.
\]
En réarrangeant l'implication,
on obtient
\[
!D(P)\land\Atom{P}{c_1,\dots,c_n}
  \Entails
!D(P')\lif\Atom{P'}{c_1,\dots,c_n}.
\]
Le théorème d'interpolation de Craig
fournit donc un !!{sentence}
$!C(c_1,\dots,c_n)$ du langage de
base enrichi de ces constantes,
ne contenant ni~$P$ ni~$P'$, tel que
\begin{align*}
!D(P)\land\Atom{P}{c_1,\dots,c_n}
  &\Entails !C(c_1,\dots,c_n),\\
!C(c_1,\dots,c_n)
  &\Entails !D(P')\lif
    \Atom{P'}{c_1,\dots,c_n}.
\end{align*}
De la première conséquence,
$!D(P)\Entails\Atom{P}{c_1,\dots,c_n}
\lif !C(c_1,\dots,c_n)$.
Renommons~$P'$ en~$P$ dans la
seconde : comme~$!C$ ne contient
aucun de ces deux symboles,
on conserve une conséquence valide,
\[
!C(c_1,\dots,c_n)\Entails
  !D(P)\lif\Atom{P}{c_1,\dots,c_n}.
\]
Les deux conséquences réunies
donnent
\[
!D(P)\Entails\bigl(
  \Atom{P}{c_1,\dots,c_n}
  \liff !C(c_1,\dots,c_n)\bigr).
\]
Les constantes $c_1,\dots,c_n$
ne figurent pas dans~$\Sigma(P)$.
Après avoir renommé au besoin les
variables liées de~$!C$ pour éviter
toute capture, remplaçons ces
constantes par $x_1,\dots,x_n$.
Puisque $\Sigma(P)\Entails !D(P)$,
la généralisation donne
\[
\Sigma(P)\Entails
  \lforall[x_1][\dots\lforall[x_n]
  [(\Atom{P}{x_1,\dots,x_n}
  \liff !C(x_1,\dots,x_n))]].
\]
Ainsi $\Sigma(P)$ définit
explicitement~$P$.
\end{proof}

\end{document}
